We study inverse calibration of a canonical ten-parameter Double Heston model
using a frozen, market-supported observation representation and a known-truth
synthetic evaluation. The production pricer is validated against an independent
adaptive-quadrature reference on 36 controlled cases. A predeclared representation
study selects a two-expiry, central-five call/put interface, denoted R2; the final
dataset contains 10{,}000 complete synthetic surfaces with stored train,
validation, and test splits.

On the untouched 1{,}250-surface test set, traditional numerical calibration is
the accuracy frontier but costs hundreds of seconds per surface. Two neural
inverse models evaluate orders of magnitude faster and always emit structurally
valid parameters, yet have materially larger parameter errors. Most importantly,
near-perfect repricing does not certify parameter recovery: traditional fits can
reprice to approximately nine significant digits while remaining materially wrong
in range-scaled parameter distance. This practical non-identifiability is retained
rather than resolved by the learned models.

A separately frozen observation-noise protocol has completed deterministic cohort
generation, strict 0\% gates, full-population neural evaluations through 1\%
noise, and the clean-gate traditional subset. Positive-noise traditional
calibration remains incomplete and is not reported as complete. A genuine PDE-
informed Model3 protocol is frozen as future active methodology; no Model3 result,
OOD result, G8 result, or positive-noise traditional completion exists in this
draft.
